\section{The Asset Manager Code}

\begin{qfdefinition}[The Asset Manager Code: a firm-level code of conduct, its purpose and benefits, the six general principles and the provisions that implement them]\label{def:asset-manager-code}
\qffield{Statement} \emph{Purpose.} The Code and Standards govern individuals; the Asset Manager Code applies to \emph{firms} that manage client assets in separate accounts or pooled funds, including hedge fund managers. It is a voluntary, global code that sets minimum ethical standards and leaves firms to write the policies that implement them, according to their size and regulators \citep[pp.~1--2]{asset-manager-code-2010}. It is meant to be adopted by senior management or the board, and to cover all employees (p.~1).

\emph{Benefits} (Interpretation, from pp.~1--2). For clients, it provides a template for selecting managers who follow sound practice and a basis for monitoring them. For firms, it signals a commitment to ethics, protects reputation, and supplies a checklist against which an existing code can be tested, which is useful to firms such as unregistered hedge funds that would otherwise have none. It also helps firms meet regulations requiring advisers to have codes of ethics \citep[p.~xi]{standards-practice-handbook-2024}.

\emph{The six general principles} (p.~3). Managers must
\begin{enumerate}
\item act in a professional and ethical manner at all times;
\item act for the benefit of clients;
\item act with independence and objectivity;
\item act with skill, competence and diligence;
\item communicate with clients in a timely and accurate manner;
\item uphold the applicable rules governing capital markets.
\end{enumerate}
\emph{Provisions} (pp.~5--7) put them into six areas: (A) loyalty to clients; (B) investment process and actions, including suitability and mandates; (C) trading, including inside information, priority of client trades, soft commissions, best execution and fair allocation; (D) risk management, compliance and support, including a compliance officer, third-party review, records, qualified staff, business continuity and firm-wide risk management; (E) performance and valuation; and (F) disclosures, a list of twelve items from conflicts to risk-management processes.

\emph{Claiming compliance.} A firm must meet every provision; partial claims are prohibited. It must state that it ``claims compliance with the CFA Institute Asset Manager Code. This claim has not been verified by CFA Institute,'' and notify CFA Institute, which does not verify or enforce compliance (p.~2).

\qffield{Practice} \emph{Mapping principles to provisions} (original): principle~2 is implemented by A.1 (client interests first), C.2 (client trades first) and C.3 (soft commissions only for client-benefiting research); principle~3 by A.3 (no gifts or relationships that impair independence) and E.2 (independent valuation); principle~5 by F.1--F.4. A firm that complies with the individual Standards through its staff may still fail the firm-level provisions: for example, it may have no business-continuity plan (D.6) or no third-party confirmation of client information (D.3).

\qffield{Conditions} Second edition (2009, reprinted 2010 and 2012, web addresses updated 2018); the guidance in its appendix is recommended, not required (p.~2).

\qffield{Pitfalls} Claiming compliance ``except for'' some provisions. Reading the claim as verified by CFA Institute.

\qffield{See also} \hyperref[meth:asset-manager-code-compliance]{assessing compliance}; \hyperref[def:code-of-ethics-standards]{individual Standards}.

\end{qfdefinition}

\begin{qfmethod}[Determining whether an asset manager's practices comply with the Asset Manager Code and recommending procedures that prevent violations: a provision-by-provision review of a hedge fund manager]\label{meth:asset-manager-code-compliance}
\qffield{Statement} Review each provision against the firm's documented practices, then recommend procedures from the Code's guidance \citep[pp.~9--22]{asset-manager-code-2010}:
\begin{itemize}
\item \emph{Loyalty}: compensation aligned with clients; a privacy policy; written gift limits of minimal value, no cash gifts, logging and disclosure (pp.~9--10).
\item \emph{Investment process}: due diligence and understanding of complex instruments, with stress tests; changes of strategy disclosed well in advance, with penalty-free redemption; an IPS for separate accounts, reviewed at least annually (pp.~10--13).
\item \emph{Trading}: information barriers; restricted and watch lists and pre-approval of IPOs and private placements for staff; soft commissions only for research used in investment decisions; best execution, with a written warning to clients who direct brokerage; pro-rata allocation of block trades and new issues (pp.~13--15).
\item \emph{Risk and compliance}: a competent compliance officer, independent where possible and reporting to the CEO or board; staff acknowledgement of the Code; third-party confirmation of client information through audits or custodian statements; records; adequate qualified staff; a tested business-continuity plan; a risk-management function independent of portfolio managers, using stress tests and scenario analysis (pp.~15--18).
\item \emph{Performance and valuation}: fair and complete performance, with backtests labelled; GIPS as the model; independent valuation of hard-to-value holdings, approved by independent directors or a third party, applied consistently (pp.~18--19).
\item \emph{Disclosures}: plain-language disclosure of conflicts, disciplinary history, process (lock-ups, leverage, derivatives), fees and costs with gross and net returns, soft commissions, quarterly performance within 30 days, valuation methods by asset class, voting and allocation policies, audit results, significant personnel changes and risk-management processes (pp.~19--22).
\end{itemize}

\qffield{Practice} \emph{Sable Ridge Capital} (original), a USD~1.2 billion credit hedge fund manager with 18 staff, wants to claim compliance. Findings:
\begin{center}\small
\begin{tabular}{@{}p{0.34\linewidth}p{0.14\linewidth}p{0.44\linewidth}@{}}
\toprule
Practice & Provision & Finding and recommended procedure \\
\midrule
CIO marks Level~3 loans monthly; performance fee on NAV & E.2, D.3 & Not consistent: conflict in self-valuation. Appoint an independent administrator or valuation agent; board-approved policy. \\
COO is also compliance officer and trader's supervisor & D.2 & Weak independence; acceptable for a small firm only if the COO reports to the board and trading reviews are done by someone else. \\
Larger investors receive monthly risk reports, others quarterly & B.3, F.1 & Consistent only if the higher service is disclosed and offered to all on the same terms. \\
Soft commissions pay for a Bloomberg terminal used by the CFO for accounting & C.3 & Not consistent: the benefit must assist investment decisions. Pay from the firm's own budget. \\
No written business-continuity plan & D.6 & Not consistent: write and test one. \\
Pitch book shows a 2015--2025 track record that includes the founder's prior firm without saying so & E.1, F.3 & Not consistent: label the predecessor record and its portability. \\
\bottomrule
\end{tabular}
\end{center}
Until all six items are fixed, Sable Ridge cannot claim compliance; a statement that it complies ``except for valuation'' is prohibited (p.~2).

\qffield{Conditions} The Code permits flexibility of implementation by firm size; the findings are judgements against its provisions.

\qffield{Pitfalls} Treating regulatory compliance as Code compliance: the Code's disclosures and independence requirements can exceed local rules. Letting portfolio managers value the assets on which their fees are based.

\qffield{See also} \hyperref[def:asset-manager-code]{Asset Manager Code}; \hyperref[meth:preventing-employer-duty-violations]{supervision}; \hyperref[meth:gips-composite-presentation]{GIPS}.

\end{qfmethod}
